\documentclass[11pt,a4paper]{article}
\usepackage[margin=2.3cm]{geometry}
\usepackage{amsmath}
\usepackage{amssymb}
\usepackage{booktabs}
\usepackage{hyperref}
\usepackage{xcolor}

\hypersetup{
  colorlinks=true,
  linkcolor=blue,
  urlcolor=blue
}

\title{SolarGen Forecast Model}
\author{SolarGen}
\date{Model version: recalibrated with measured history from 2026-05-04 through 2026-05-12}

\begin{document}
\maketitle

\section{Purpose and Scope}

SolarGen estimates hourly rooftop photovoltaic generation for a 10 kWp south-facing rooftop array in OHZ / Osterholz-Scharmbeck. The model is used by both the browser forecast and the local forecast-history database.

The model has two calibration layers:

\begin{itemize}
  \item a clear-sky anchor day: 50.23 kWh measured on 2026-05-01;
  \item a weather-response recalibration period: measured forecast history from 2026-05-04 through 2026-05-12.
\end{itemize}

The hourly input data comes from Open-Meteo:

\[
I_h,\ C_h,\ R_h,\ T_h
\]

where \(I_h\) is global tilted irradiance in \(\mathrm{W/m^2}\), \(C_h\) is cloud cover in percent, \(R_h\) is precipitation in mm, and \(T_h\) is ambient temperature in \(^{\circ}\mathrm{C}\). Daily inputs are precipitation sum, mean cloud cover, and sunshine duration.

\section{Notation}

\begin{center}
\begin{tabular}{ll}
\toprule
Symbol & Meaning \\
\midrule
\(h\) & local hour, 0 through 23 \\
\(P\) & configured PV capacity in kWp, default \(10\) \\
\(F\) & feed-in cap in kW, default \(6\) \\
\(I_h\) & Open-Meteo global tilted irradiance, \(\mathrm{W/m^2}\) \\
\(C_h\) & hourly cloud cover fraction, \(C_h = \mathrm{cloud\_pct}_h / 100\) \\
\(R_h\) & hourly precipitation, mm \\
\(T_h\) & ambient air temperature, \(^{\circ}\mathrm{C}\) \\
\(S\) & calibration scale solved from the clear-sky anchor day \\
\bottomrule
\end{tabular}
\end{center}

All hourly kW averages are treated as kWh over that hour.

\section{Weather-Adjusted Irradiance}

\subsection{Low-Irradiance Cloud Uplift}

Very dark overcast days produced more rooftop output than raw tilted irradiance implied. The model therefore boosts low-irradiance cloudy hours:

\[
L_h = \left[\operatorname{clamp}\left(1 - \frac{I_h}{577}, 0, 1\right)\right]^{1.855}
\]

\[
M_h = \min\left(1.748,\ 1 + 2.2 C_h L_h\right)
\]

Clear hours keep \(M_h = 1\). The boost is strongest when cloud cover is high and irradiance is low.

\subsection{Daily Weather Damping}

The May 12 actual showed the opposite failure mode: bright hourly irradiance under wet, low-sun, high-cloud daily conditions can over-forecast. The model computes:

\[
D_{\mathrm{rain}} = 0.195 \cdot \operatorname{smoothstep}(0,\ 3.116,\ R_{\mathrm{day}})
\]

\[
D_{\mathrm{sun}} = 0.747 \cdot \operatorname{smoothstep}(10.267,\ 3.755,\ H_{\mathrm{sun}})
\]

\[
D_{\mathrm{cloud}} = 0.458 \cdot \operatorname{smoothstep}(61.016,\ 93.522,\ C_{\mathrm{day}})
\]

\[
D_{\mathrm{day}} =
\operatorname{clamp}\left(D_{\mathrm{rain}} + D_{\mathrm{sun}} + D_{\mathrm{cloud}}, 0, 1.2\right)
\]

where \(R_{\mathrm{day}}\) is daily precipitation in mm, \(H_{\mathrm{sun}}\) is sunshine duration in hours, and \(C_{\mathrm{day}}\) is mean daily cloud cover in percent.

\subsection{Bright High-Cloud Damping}

The bright-hour weight is:

\[
B_h = \operatorname{smoothstep}(265,\ 818,\ I_h)
\]

The bright-cloud damping factor is:

\[
Q_h =
\operatorname{clamp}\left(1 - 0.763 C_h B_h D_{\mathrm{day}},\ 0.25,\ 1.1\right)
\]

\subsection{Hourly Rain Damping}

\[
U_h = \frac{1}{1 + 0.955 R_h}
\]

\subsection{Adjusted Irradiance}

The weather-adjusted irradiance is:

\[
I'_h = I_h \cdot M_h \cdot Q_h \cdot U_h
\]

\section{Temperature Derating}

Panel temperature is approximated from adjusted irradiance and ambient temperature:

\[
T_{\mathrm{panel},h} = T_h + \frac{I'_h}{800} \cdot 20
\]

The derating factor is:

\[
A_h =
\operatorname{clamp}\left(1 - 0.0035(T_{\mathrm{panel},h} - 25),\ 0.82,\ 1.06\right)
\]

\section{Theoretical PV Output}

Before applying the site-specific rooftop profile:

\[
G^{\mathrm{theory}}_h =
\frac{I'_h}{1000} \cdot P \cdot S \cdot A_h
\]

The calibration scale \(S\) is solved numerically so that the clear-sky irradiance curve for 2026-05-01 produces exactly 50.23 kWh after the rooftop profile is applied.

\section{Rooftop Profile}

The rooftop profile maps theoretical generation to the observed installation behavior. It models morning limitation, the main production window, evening drop-off, and diffuse-light behavior under cloud.

\subsection{Clear-Sky Branch}

The low-output branch is:

\[
G^{\mathrm{low}}_h =
\min\left(0.286 G^{\mathrm{theory}}_h,\ 1.44\right)
\]

The clear/direct branch is:

\[
G^{\mathrm{sun}}_h =
\begin{cases}
G^{\mathrm{low}}_h, & h < 10.67 \\
G^{\mathrm{theory}}_h, & 10.67 \leq h < 16.72 \\
(1 - p_h)G^{\mathrm{theory}}_h + p_h G^{\mathrm{low}}_h, & 16.72 \leq h < 17.61 \\
G^{\mathrm{low}}_h, & h \geq 17.61
\end{cases}
\]

where

\[
p_h = \frac{h - (17.61 - 0.89)}{0.89}
\]

\subsection{Cloud/Diffuse Branch}

Clouds increase the diffuse-light share:

\[
W_h = \operatorname{smoothstep}(57.94,\ 93.47,\ 100 C_h)
\]

The diffuse branch uses the hour center \(x_h = h + 0.5\):

\[
G^{\mathrm{diff}}_h =
G^{\mathrm{theory}}_h
\cdot \operatorname{smoothstep}(4.5,\ 7.9,\ x_h)
\cdot \left[1 - \operatorname{smoothstep}(13.18,\ 19.9,\ x_h)\right]
\cdot 1.2
\]

\subsection{Rooftop Generation Before Curtailment}

\[
G_h =
(1 - W_h)G^{\mathrm{sun}}_h + W_hG^{\mathrm{diff}}_h
\]

For hours 5 and 6 only, a small dawn diffuse floor is applied:

\[
G_h = \max(G_h,\ 0.2 C_h \cdot P / 10)
\]

\section{Curtailment and Delivered PV}

The feed-in cap is applied to hourly rooftop generation:

\[
G^{\mathrm{curt}}_h = \max(0,\ G_h - F)
\]

\[
G^{\mathrm{deliv}}_h = G_h - G^{\mathrm{curt}}_h
\]

Daily rooftop generation is:

\[
G_{\mathrm{day}} = \sum_{h=0}^{23} G_h
\]

\section{Household Load and Battery Accounting}

The default household load is:

\[
L^{\mathrm{home}}_h =
L_{\mathrm{base}}
+ \mathbf{1}_{8 \leq h < 18}L_{\mathrm{day}}
+ \mathbf{1}_{18 \leq h < 23}L_{\mathrm{eve}}
+ \mathbf{1}_{6 \leq h < 8}(0.45L_{\mathrm{day}})
\]

PV first covers same-hour load. Remaining load is covered from the battery if available. PV surplus charges the battery with 94 percent charge efficiency. Remaining surplus is exported.

\section{Calibration Procedure}

Calibration is split into a structural clear-sky fit and a weather-response fit. This keeps the model identifiable: the clear-sky day fixes the site scale and rooftop geometry, while the historical cloudy and rainy days tune how meteorology changes production.

\subsection{Clear-Sky Scale}

Let \(A\) be the anchor day, 2026-05-01, with measured production

\[
Y_A = 50.23\ \mathrm{kWh}.
\]

For any trial scale \(S\), the model produces an anchor-day total

\[
\widehat{Y}_A(S) = \sum_{h=0}^{23} G_h(S).
\]

The scale is solved numerically as the one-dimensional root

\[
S^{\star}
= \arg\min_S \left(\widehat{Y}_A(S) - Y_A\right)^2,
\qquad
\widehat{Y}_A(S^{\star}) = Y_A.
\]

This step fixes the absolute output level for the 10 kWp array after temperature derating and the rooftop profile.

\subsection{Weather-Response Fit}

The weather-response parameters are fitted against measured daily totals from 2026-05-04 through 2026-05-12. For each measured day \(d\), let \(Y_d\) be actual generation and \(\widehat{Y}_d(\theta)\) be the modeled total under parameter vector \(\theta\):

\[
\theta =
\left(
\theta_{\mathrm{cloud}},
\theta_{\mathrm{rain}},
\theta_{\mathrm{sun}},
\theta_{\mathrm{profile}}
\right).
\]

The fitting objective minimizes daily error while keeping the anchor scale fixed:

\[
\theta^{\star}
= \arg\min_{\theta}
\sum_{d \in \mathcal{D}}
w_d
\left(\widehat{Y}_d(\theta) - Y_d\right)^2,
\qquad
\mathcal{D} = \{2026\text{-}05\text{-}04,\ldots,2026\text{-}05\text{-}12\}.
\]

Days with incomplete actuals are excluded from \(\mathcal{D}\). The calibrated parameters primarily control four effects:

\begin{itemize}
  \item low-irradiance cloud uplift, \(M_h\), which raises modeled output on dark diffuse days;
  \item daily weather damping, \(D_{\mathrm{day}}\), which suppresses over-optimistic bright-hour forecasts on wet, low-sun, high-cloud days;
  \item hourly rain damping, \(U_h\), which applies immediate precipitation loss;
  \item rooftop timing parameters, which encode the observed morning limitation and evening drop-off.
\end{itemize}

\subsection{Model Selection Criterion}

The selected parameter set is the simplest candidate that reproduces the historical totals without removing the physical structure of the model. The primary score is daily mean absolute error:

\[
\operatorname{MAE}
=
\frac{1}{|\mathcal{D}|}
\sum_{d \in \mathcal{D}}
\left|\widehat{Y}_d - Y_d\right|.
\]

Relative accuracy is tracked with daily mean absolute percentage error:

\[
\operatorname{MAPE}
=
\frac{100}{|\mathcal{D}|}
\sum_{d \in \mathcal{D}}
\left|
\frac{\widehat{Y}_d - Y_d}{Y_d}
\right|.
\]

The calibration period must be updated in this document and in the app header whenever the weather-response fit is regenerated.

\section{Calibration Quality}

Replaying the measured days from 2026-05-04 through 2026-05-12 gives:

\begin{center}
\begin{tabular}{lr}
\toprule
Metric & Value \\
\midrule
Days compared & 7 \\
Actual total & 235.33 kWh \\
Model total & 238.13 kWh \\
Daily mean absolute error & 0.44 kWh \\
Daily RMSE & 0.55 kWh \\
Daily MAPE & 1.35 percent \\
Worst daily absolute percentage error & 2.43 percent \\
\bottomrule
\end{tabular}
\end{center}

These are in-sample statistics because the measured days were used for recalibration. Future actuals should be used as out-of-sample checks.

\section{Implementation References}

The corresponding implementation lives in:

\begin{itemize}
  \item \texttt{src/config.js}: constants and calibration parameters;
  \item \texttt{src/model.js}: weather adjustment, rooftop profile, storage, export, and totals;
  \item \texttt{tests/model.test.mjs}: calibration replay tests and model invariants.
\end{itemize}

\end{document}
